\section{Từ nguyên tử đến tích phân}
\label{sec:ch06-tu-nguyen-tu-den-tich-phan}

\index{độ đo|(}%
\index{tích phân}%
\tn{Hàm chỉ báo}{indicator function} giải thích một vùng, trong khi mô hình thực
tế thay đổi trên nhiều vùng. Bài báo đi từ vùng tới mô hình qua bốn bước, và
hình~\ref{fig:ch06-xay-attribution} vẽ đúng bốn bước ấy theo chiều xây
dựng~\autocite{p17firstprinciples}. Ở đây bài báo giới hạn miền vào khối đơn vị
$[0,1]^d$, với $d$ là số chiều của không gian đầu vào, và giả sử $f$ liên tục
trên khối ấy; hai giới hạn này là điều kiện để bước lấy giới hạn phía sau dùng
được.

Bước thứ nhất là \tn{quy tắc nguyên tử}{atomic attribution rule} của mục
trước. Với mỗi hình hộp $R$, quy tắc gán cho đặc trưng $j$ tại $x$ một số, và
ta ký hiệu số ấy là $\mu_{j,x}(R)$. Ký hiệu này chỉ ghi lại quy tắc, chưa thêm
gì mới.

Bước thứ hai xấp xỉ $f$. Chia $[0,1]^d$ thành một lưới các hình hộp nhỏ
$R_1,\dots,R_n$, lấy giá trị của $f$ tại một đỉnh $y_k$ của mỗi hộp, rồi đặt
$f_n=\sum_k f(y_k)\,\mathbf{1}_{R_k}$. Hàm $f_n$ là hàm bậc thang, hằng trên
từng hộp, và vì $f$ liên tục nên lưới càng mịn thì $f_n$ càng sát $f$ đều trên
toàn khối.

Bước thứ ba dùng tính tuyến tính, tiên đề duy nhất mà chương này giữ lại.
Attribution của một tổng bằng tổng các attribution, nên
\begin{equation}
\label{eq:ch06-tong-riemann}
\varphi_j(x,f_n)=\sum_{k} f(y_k)\,\mu_{j,x}(R_k).
\end{equation}
Vế phải là một tổng Riemann: giá trị hàm tại mỗi hộp nhân với trọng lượng của
hộp ấy. Chỗ khác với tổng Riemann quen thuộc là trọng lượng không phải thể tích
của hộp mà là số $\mu_{j,x}(R_k)$ do quy tắc nguyên tử gán.

Bước thứ tư lấy giới hạn khi lưới mịn dần, và để giới hạn ấy có nghĩa, bài báo
cần hai điều. Thứ nhất, các số $\mu_{j,x}(R)$ mà quy tắc nguyên tử gán cho các
hình hộp phải là giá trị của một
\tn{độ đo Borel hữu hạn có dấu}{finite signed Borel measure} trên $[0,1]^d$,
tức một cách gán trọng lượng, có thể âm, cho mọi tập Borel của khối, là mọi
vùng dựng được từ các hình hộp bằng hợp, giao và phần bù đếm được, không chỉ
cho các hình hộp. Bài báo lấy điều này làm giả thiết: mỗi quy tắc nguyên tử
được cho sẵn dưới dạng một họ độ đo như thế. Thứ hai, attribution phải ổn định trước hội
tụ đều của mô hình: nếu $f_n$ tiến tới $f$ đều thì $\varphi_j(x,f_n)$ phải
tiến tới $\varphi_j(x,f)$. Không có điều này, giới hạn của các tổng trong
phương trình~\ref{eq:ch06-tong-riemann} là một con số không nói gì về $f$. Với
hai điều ấy, tổng Riemann hội tụ về tích phân của $f$ theo
\tn{độ đo}{measure} $\mu_{j,x}$, và attribution của $f$ là

\begin{equation}
\label{eq:ch06-attribution-do-do}
\varphi_j(x,f)=\int_{[0,1]^d} f(y)\,\mathrm{d}\mu_{j,x}(y).
\end{equation}

Bài báo còn nêu chiều ngược lại, rút từ định lý biểu diễn Riesz của giải tích
hàm: mọi attribution tuyến tính và ổn định trước hội tụ đều, xét trên các hàm
liên tục của $[0,1]^d$, đều có dạng của phương
trình~\ref{eq:ch06-attribution-do-do} với một họ độ đo duy
nhất~\autocite{p17firstprinciples}. Vậy chọn quy tắc nguyên tử và chọn độ đo là
cùng một việc, nhìn từ hai phía.

Ví dụ hai chiều của mục~\ref{sec:ch06-ham-chi-bao} cho thấy một độ đo trông ra
sao. Quy tắc thứ hai ở đó gán cho đặc trưng thứ nhất số $b_2-a_2$ khi
$x_1\in(a_1,b_1]$, và số ấy đúng là trọng lượng mà độ đo sau đây gán cho hình
hộp $R$: toàn bộ trọng lượng tập trung trên đường thẳng $y_1=x_1$ và trải đều
dọc theo tọa độ thứ hai. Tích phân của $f$ theo độ đo ấy là
$\int_0^1 f(x_1,y_2)\,\mathrm{d}y_2$: giữ tọa độ đang xét tại giá trị của $x$,
lấy trung bình tọa độ kia trên cả khoảng.

Ở đây $\mu_{j,x}$ gắn với đặc trưng $j$ và đầu vào $x$, và nó không phải một
chi tiết cài đặt. Nó quyết định những điểm $y$ nào của không gian đầu vào được
tính, và chúng được tính với dấu cùng trọng số ra sao.

\begin{figure}[htbp]
  \centering
  \begin{tikzpicture}[
  box/.style={draw=black!65, rounded corners=2pt, align=center,
    minimum width=3.4cm, minimum height=1.1cm, fill=black!5},
  data/.style={draw=black!65, rounded corners=2pt, align=center,
    minimum width=3.4cm, minimum height=1.1cm, fill=black!15},
  arrow/.style={-{Stealth[length=2mm]}, draw=black!65, thick}
]
  \node[box] (atom) {Quy tắc attribution\\cho hàm chỉ báo $\mathbf{1}_R$};
  \node[data, right=11mm of atom] (measure) {Độ đo\\$\mu_{j,x}$};
  \node[box, right=11mm of measure] (sum) {Xấp xỉ $f$ bằng tổng\\trên các vùng $R$};
  \node[data, below=11mm of sum] (integral) {Giới hạn:\\$\varphi_j(x,f)=\int f\,\mathrm{d}\mu_{j,x}$};

  \draw[arrow] (atom) -- (measure);
  \draw[arrow] (measure) -- (sum);
  \draw[arrow] (sum) -- (integral);
\end{tikzpicture}

  \caption{Bốn bước xây attribution của Taimeskhanov và Garreau, đúng thứ tự
  của mục này: quy tắc nguyên tử cho hàm chỉ báo, phép xấp xỉ $f$ bằng tổng
  trên các hình hộp, tính tuyến tính cho tổng Riemann của
  phương trình~\ref{eq:ch06-tong-riemann}, rồi giới hạn đưa attribution về
  tích phân trong phương trình~\ref{eq:ch06-attribution-do-do}. Sơ đồ là cách
  sách này đọc lại cơ chế của bài báo~\autocite{p17firstprinciples}.}
  \label{fig:ch06-xay-attribution}
\end{figure}
